% Moved verbatim from the main text during the length revision
% (was common_body.tex lines 407-413). Content unchanged.
The distinction this study draws between the presence of information in a representation and a specific estimator's ability to use it is closely related to, and constrained by, the probing-classifier literature's own caution that a supervised probe's success does not establish that a model itself uses the decoded information in the way the probe accesses it \citep{hewitt2019designing, belinkov2022probing}. This study makes a narrower claim than that literature warns against overreaching toward: the probes here are not offered as evidence about what the disentanglement objective computes internally, only as an independent measurement of what remains extractable from its output representation, without characterizing the mechanism by which that representation was produced.

This finding also bears on the geometry-performance association reported across frozen foundation-model backbones by \citet{janiak2026geometry} (Section~\ref{sec:rw-geometry}). That association is measured between representations whose provenance -- pretraining corpus, objective, architecture, and scale -- differs simultaneously, so it cannot by itself distinguish a geometric explanation from a shared dependence on pretraining. The present study varies orthogonality strength within a fixed architecture and representation dimensionality and finds no corresponding geometry--AUROC association for the measured quantities tested here, including condition number (Section~\ref{sec:geometry-auroc}). Whether this absence of association reflects a property specific to domain-adversarial disentanglement, to this dataset pair, or more generally to interventions that hold pretraining fixed, is a question this study's design cannot resolve on its own and that we regard as a natural target for direct follow-up.

More broadly, these findings suggest that evaluating learned representations only through downstream predictive performance may overlook how a downstream reliability score orders shifted inputs. Raw geometric summaries, scorer-relevant geometry, operational score direction, orientation-free scalar separability, and supervised decodability are related but distinct properties; change in one should not automatically be interpreted as change in the others.

For practitioners combining domain-adversarial representation learning with a distance-based reliability gate, this audit implies a specific caution: a score can weakly distinguish the targeted domain yet assign that distinction the wrong operational polarity. Validation should therefore test both discrimination and direction on calibration data independent of final evaluation, rather than infer operational validity from an orientation-free metric or from general model qualification.
